% !TEX root = exercices_s3.tex
% ============================================================================
% HEC Montréal MBA -- ECON50803: Macroeconomic Environment
% Shared Preamble for Exercise Handouts
% ============================================================================

% --- Document class ---------------------------------------------------------
\documentclass[11pt, letterpaper]{article}
\usepackage[margin=2.5cm]{geometry}

% --- Fonts ------------------------------------------------------------------
\usepackage[sfdefault,light]{FiraSans}
\usepackage{FiraMono}
\usepackage[T1]{fontenc}

% --- Colors (same as slides) ------------------------------------------------
\usepackage{xcolor}
\definecolor{HECnavy}{HTML}{002855}
\definecolor{HECgreen}{HTML}{26d07c}
\definecolor{HECcoral}{HTML}{ff585d}
\definecolor{HECtext}{HTML}{2D3748}
\definecolor{HECgray}{HTML}{94A3B8}
\definecolor{HEClightgray}{HTML}{F1F5F9}
\definecolor{HEClightblue}{HTML}{E8EDF3}

% --- Packages ---------------------------------------------------------------
\usepackage{amsmath, graphicx, tikz, booktabs, tabularx, hyperref}
\usepackage[inline]{enumitem}
\usepackage{etoolbox}
\usepackage{comment}
\usepackage{needspace}
\usepackage{tcolorbox}
\tcbuselibrary{skins, breakable}
\usepackage{titlesec}
\usepackage{fancyhdr}
\usepackage{lastpage}

% --- Paths ------------------------------------------------------------------
\graphicspath{{../Slides/Figures/}}

% --- Hyperlinks -------------------------------------------------------------
\hypersetup{
   colorlinks,
   linkcolor=HECnavy,
   urlcolor=HECnavy,
   citecolor=HECtext
}

% --- Answer toggle ----------------------------------------------------------
\newbool{answers}
\ifdefined\showanswers
   \booltrue{answers}
\else
   \boolfalse{answers}
\fi

% Answer environment: shown as green box when answers=true, excluded otherwise
\ifbool{answers}{%
   \newtcolorbox{reponse}[1][Réponse]{%
      colback=HECgreen!5,
      colframe=HECgreen,
      boxrule=0pt,
      leftrule=3pt,
      arc=4pt,
      title={\textcolor{HECgreen}{\bfseries #1}},
      coltitle=HECgreen,
      fonttitle=\bfseries,
      attach title to upper={\\[4pt]},
      left=8pt, right=8pt, top=4pt, bottom=4pt,
      before={\par\nopagebreak\medskip\nopagebreak}
   }%
}{%
   \excludecomment{reponse}%
}

% --- Custom box environments (adapted from slides) --------------------------

% Key Insight (green accent)
\newtcolorbox{keyinsight}[1][Point clé]{%
   colback=HECgreen!5,
   colframe=HECgreen,
   boxrule=0pt,
   leftrule=3pt,
   arc=4pt,
   title={\textcolor{HECgreen}{\bfseries #1}},
   coltitle=HECgreen,
   fonttitle=\bfseries,
   attach title to upper={\\[4pt]},
   left=8pt, right=8pt, top=4pt, bottom=4pt,
   breakable
}

% Business Implication (navy accent)
\newtcolorbox{bizimplication}[1][Implication d'affaires]{%
   colback=HECnavy!5,
   colframe=HECnavy,
   boxrule=0pt,
   leftrule=3pt,
   arc=4pt,
   title={\textcolor{HECnavy}{\bfseries #1}},
   coltitle=HECnavy,
   fonttitle=\bfseries,
   attach title to upper={\\[4pt]},
   left=8pt, right=8pt, top=4pt, bottom=4pt,
   breakable
}

% Warning / Important (coral accent)
\newtcolorbox{warning}[1][Important]{%
   colback=HECcoral!5,
   colframe=HECcoral,
   boxrule=0pt,
   leftrule=3pt,
   arc=4pt,
   title={\textcolor{HECcoral}{\bfseries #1}},
   coltitle=HECcoral,
   fonttitle=\bfseries,
   attach title to upper={\\[4pt]},
   left=8pt, right=8pt, top=4pt, bottom=4pt,
   breakable
}

% Definition (gray background)
\newtcolorbox{defbox}[1][Définition]{%
   colback=HEClightgray,
   colframe=HEClightgray,
   boxrule=0pt,
   leftrule=3pt,
   arc=4pt,
   left=8pt, right=8pt, top=4pt, bottom=4pt,
   title={\textcolor{HECnavy}{\bfseries #1}},
   coltitle=HECnavy,
   fonttitle=\bfseries,
   attach title to upper={\\[4pt]},
   breakable
}

% Exercise (coral accent)
\newtcolorbox{exercisebox}[1][Exercice]{%
   colback=HECcoral!5,
   colframe=HECcoral,
   boxrule=0pt,
   leftrule=3pt,
   arc=4pt,
   title={\textcolor{HECcoral}{\bfseries #1}},
   coltitle=HECcoral,
   fonttitle=\bfseries,
   attach title to upper={\\[4pt]},
   left=8pt, right=8pt, top=4pt, bottom=4pt,
   breakable
}

% --- Section formatting (titlesec) ------------------------------------------
\titleformat{\section}
   {\clearpage\Large\bfseries\color{HECnavy}}
   {\thesection.}{0.5em}{}
   [\vspace{2pt}{\color{HECgreen}\rule{2.5cm}{1.5pt}}]
\titlespacing*{\section}{0pt}{1.5em}{0.8em}

\titleformat{\subsection}
   {\large\bfseries\color{HECnavy}}
   {\thesubsection.}{0.5em}{}
\titlespacing*{\subsection}{0pt}{1em}{0.5em}

% --- Header/Footer (fancyhdr) ----------------------------------------------
\pagestyle{fancy}
\fancyhf{}
\renewcommand{\headrulewidth}{0pt}
\renewcommand{\footrulewidth}{0.4pt}
\renewcommand{\footrule}{\hfill{\color{HECgray}\rule{0.9\textwidth}{0.4pt}}\hfill\vspace{4pt}}
\fancyfoot[L]{\small\textcolor{HECgray}{ECON 50803\enspace\textbar\enspace Environnement macroéconomique}}
\fancyfoot[R]{\small\textcolor{HECgray}{\thepage\enspace/\enspace\pageref{LastPage}}}

% --- Enumitem defaults ------------------------------------------------------
\setlist[itemize]{leftmargin=1.5em, itemsep=3pt}
\setlist[enumerate]{leftmargin=1.5em, itemsep=3pt}

% --- MC question helpers ----------------------------------------------------
\newcounter{qcmcounter}
\newcommand{\qcm}[1]{%
   \stepcounter{qcmcounter}%
   \par\vspace{1.2em}%
   \needspace{12\baselineskip}%
   \noindent\textbf{\textcolor{HECnavy}{\theqcmcounter.}}\enspace #1%
   \vspace{0.3em}%
   \nopagebreak[4]%
}

\newcommand{\choix}[1]{%
   \nopagebreak[4]%
   \begin{enumerate}[label=\textcolor{HECgray}{(\alph*)}, leftmargin=2em, itemsep=1pt, topsep=2pt]
      #1
   \end{enumerate}
}

% --- Short answer counter ---------------------------------------------------
\newcounter{shortcounter}
\newcommand{\shortq}[1]{%
   \stepcounter{shortcounter}%
   \par\vspace{1.5em}%
   \needspace{4\baselineskip}%
   \noindent\textbf{\textcolor{HECnavy}{Question \theshortcounter.}}\enspace #1%
   \vspace{0.3em}%
}

% --- VFI (Vrai/Faux/Incertain) helper --------------------------------------
\newcommand{\vfi}[1]{%
   \stepcounter{shortcounter}%
   \par\vspace{1.5em}%
   \needspace{10\baselineskip}%
   \noindent\textbf{\textcolor{HECnavy}{Question \theshortcounter.}}\enspace\textit{Vrai, faux ou incertain?} #1%
   \vspace{0.3em}%
}

% --- Title page command -----------------------------------------------------
\newcommand{\exercisetitlepage}[2]{%
   % #1 = session number, #2 = session title
   \thispagestyle{empty}
   \begin{tikzpicture}[remember picture, overlay]
      % Navy sidebar
      \fill[HECnavy] (current page.north west) rectangle
         ([xshift=0.8cm]current page.south west);
      % Green accent line
      \fill[HECgreen] ([xshift=0.8cm]current page.north west) rectangle
         ([xshift=0.85cm]current page.south west);
   \end{tikzpicture}
   \vspace*{2cm}
   \hspace{0.8cm}%
   \begin{minipage}{0.85\textwidth}
      {\fontsize{28}{34}\selectfont\bfseries\textcolor{HECnavy}{Exercices --- Séance #1}}
      \\[14pt]
      {\Large\textcolor{HECnavy}{#2}}
      \\[14pt]
      {\color{HECgreen}\rule{3cm}{2pt}}
      \\[14pt]
      {\large\textcolor{HECgray}{ECON 50803 --- Environnement macroéconomique}}
      \\[8pt]
      {\large\textcolor{HECgray}{Jean-Félix Brouillette}}
      \\[20pt]
      \ifbool{answers}{%
         {\Large\bfseries\textcolor{HECgreen}{AVEC SOLUTIONS}}%
      }{%
         {\Large\bfseries\textcolor{HECgray}{SANS SOLUTIONS}}%
      }
   \end{minipage}
   \vfill
   \hspace{0.8cm}%
   {\small\textcolor{HECgray}{HEC Montr\'eal\enspace\textbar\enspace MBA}}
   \vspace{1cm}
   \newpage
}

% --- Mathematical commands (from slides) ------------------------------------
\newcommand{\pctchg}[1]{\%\Delta #1}


\begin{document}

\exercisetitlepage{3}{Marché du travail, inégalités et IA}

% ============================================================================
\section{Questions à choix multiples}
% ============================================================================

\setcounter{qcmcounter}{0}

% ---------- Indicateurs du marché du travail (Q1--Q8) ----------

\qcm{Un retraité de 67 ans qui ne cherche pas d'emploi est classé comme :}
\choix{
   \item Employé
   \item Chômeur
   \item Inactif
   \item Hors de la population en âge de travailler
}

\begin{reponse}
\textbf{(c)} La population en âge de travailler (15 ans et plus) se divise en trois catégories mutuellement exclusives : \textbf{employés} ($E$), \textbf{chômeurs} ($U$, sans emploi mais en recherche active) et \textbf{inactifs} ($I$, ne cherchent pas d'emploi). Un retraité qui ne cherche pas d'emploi est \textbf{inactif}. Il fait partie de la population en âge de travailler, mais pas de la \textbf{population active} ($E + U$). (d) serait correct uniquement s'il avait moins de 15 ans.
\end{reponse}

\qcm{Si des travailleurs découragés cessent de chercher un emploi et quittent la population active, l'effet immédiat est :}
\choix{
   \item Le taux de chômage augmente
   \item Le taux de chômage diminue
   \item Le taux de participation augmente
   \item Le taux d'emploi augmente
}

\begin{reponse}
\textbf{(b)} Les travailleurs découragés passent de la catégorie \guillemotleft~chômeurs~\guillemotright\ ($U$) à la catégorie \guillemotleft~inactifs~\guillemotright\ ($I$). Le numérateur du taux de chômage ($U$) diminue, et le dénominateur ($E + U$) diminue aussi. Puisque le taux de chômage est inférieur à 100\,\% (le numérateur est plus petit que le dénominateur), retirer une unité des deux fait \textbf{diminuer} le ratio. C'est pourquoi le taux de chômage seul peut être trompeur : une baisse peut refléter des retraits du marché plutôt que des embauches. Le taux de participation ($= (E+U)/\text{Pop.}$) diminue aussi.
\end{reponse}

\qcm{Dans une économie, la population en âge de travailler est de 200\,000 personnes. La population active est de 130\,000 et il y a 117\,000 employés. Le taux de participation est de :}
\choix{
   \item 58.5\,\%
   \item 90.0\,\%
   \item 65.0\,\%
   \item 10.0\,\%
}

\begin{reponse}
\textbf{(c)} Le taux de participation $= \frac{E + U}{\text{Pop.}} = \frac{130\,000}{200\,000} = 65.0\,\%$. Les distracteurs correspondent à d'autres ratios : (a)~$E/\text{Pop.} = 117/200 = 58.5\,\%$ (le taux d'emploi, pas le taux de participation); (b)~$E/(E+U) = 117/130 = 90.0\,\%$ (la part des employés parmi les actifs, pas un indicateur standard); (d)~$U/(E+U) = 13/130 = 10.0\,\%$ (le taux de chômage). Le taux de participation mesure la proportion de la population en âge de travailler qui est dans la \textbf{population active} (employés + chômeurs).
\end{reponse}

\qcm{Dans une économie, la population en âge de travailler est de 500\,000 personnes. Il y a 350\,000 employés, 50\,000 chômeurs et 100\,000 inactifs. Le taux de chômage est de :}
\choix{
   \item 10.0\,\%
   \item 12.5\,\%
   \item 14.3\,\%
   \item 20.0\,\%
}

\begin{reponse}
\textbf{(b)} Le taux de chômage $= \frac{U}{E + U} = \frac{50\,000}{350\,000 + 50\,000} = \frac{50\,000}{400\,000} = 12.5\,\%$. Attention : le dénominateur est la \textbf{population active} ($E + U = 400\,000$), pas la population en âge de travailler (500\,000). Les distracteurs correspondent à d'autres erreurs de calcul : (a)~$U/\text{Pop.} = 50/500 = 10\,\%$ (divise par la population totale au lieu de la population active); (c)~$U/E = 50/350 \approx 14.3\,\%$ (divise par les employés seulement); (d)~$I/\text{Pop.} = 100/500 = 20\,\%$ (la part des inactifs, pas le taux de chômage).
\end{reponse}

\qcm{Le vieillissement de la population tend à :}
\choix{
   \item Augmenter le taux de participation
   \item Réduire le taux de participation
   \item Augmenter le taux de chômage
   \item N'avoir aucun effet sur les indicateurs du marché du travail
}

\begin{reponse}
\textbf{(b)} Le vieillissement de la population augmente la proportion de personnes âgées (retraités) dans la population en âge de travailler. Les retraités sont \textbf{inactifs} (ne cherchent pas d'emploi), ce qui réduit le ratio population active / population en âge de travailler, soit le \textbf{taux de participation}. C'est un défi majeur pour les économies avancées : moins de travailleurs par rapport à la population totale $\to$ pression sur les systèmes de retraite et la croissance du PIB par habitant.
\end{reponse}

\qcm{Pendant la pandémie de COVID-19, de nombreux travailleurs ont quitté la population active (retraite anticipée, retrait volontaire), tandis que des licenciements massifs ont eu lieu. L'effet combiné sur le taux de chômage est :}
\choix{
   \item Nécessairement une hausse, car les licenciements dominent
   \item Nécessairement une baisse, car les retraits dominent
   \item Ambigu, car les deux effets vont en sens contraire
   \item Nul, car les deux effets s'annulent parfaitement
}

\begin{reponse}
\textbf{(c)} Trois effets agissent simultanément sur le taux de chômage ($U/(E+U)$) :
\begin{itemize}
   \item \textbf{Licenciements :} $E \downarrow$, $U \uparrow$ $\to$ le taux de chômage \textbf{augmente}
   \item \textbf{Retraites anticipées :} des employés quittent directement la population active ($E \downarrow$, $I \uparrow$). Le numérateur ($U$) ne change pas, mais le dénominateur ($E + U$) diminue $\to$ le taux de chômage \textbf{augmente aussi}
   \item \textbf{Retraits découragés :} des chômeurs cessent de chercher ($U \downarrow$, $I \uparrow$). Le numérateur et le dénominateur diminuent; puisque $U < E + U$, le ratio \textbf{baisse}
\end{itemize}
Les deux premiers effets poussent le taux de chômage à la hausse, le troisième à la baisse. L'effet net dépend de l'ampleur relative de chacun. C'est pourquoi le taux de chômage seul est un indicateur incomplet --- il faut aussi regarder le taux de participation et le taux d'emploi.
\end{reponse}

\qcm{Entre janvier et décembre, le taux de chômage d'un pays passe de 10\,\% à 6\,\% et le taux de participation passe de 60\,\% à 65\,\%. On peut conclure que :}
\choix{
   \item La baisse du chômage est trompeuse, car elle est due au retrait de travailleurs découragés
   \item Le marché du travail s'est détérioré
   \item La population en âge de travailler a diminué
   \item Le marché du travail s'est véritablement amélioré, avec à la fois plus d'emplois et plus de travailleurs actifs
}

\begin{reponse}
\textbf{(d)} La baisse du taux de chômage pourrait être trompeuse si elle était accompagnée d'une baisse du taux de participation (signe de travailleurs découragés). Mais ici, le taux de participation a \textbf{augmenté} de 60\,\% à 65\,\%, ce qui signifie que davantage de personnes sont entrées dans la population active. La baisse du chômage reflète donc une \textbf{véritable amélioration} : plus de gens cherchent du travail \textbf{et} en trouvent. C'est l'exemple parfait de pourquoi il faut regarder les trois indicateurs conjointement.
\end{reponse}

\qcm{Alors qu'une profonde récession s'installait en Espagne après 2009, un grand nombre de résidents adultes chômeurs ont émigré vers d'autres pays. L'effet de cette émigration \textbf{seule} est :}
\choix{
   \item Une baisse du taux d'activité (participation)
   \item Une hausse du taux de participation
   \item Une hausse du taux de chômage
   \item Aucun effet sur les indicateurs du marché du travail
}

\begin{reponse}
\textbf{(a)} Les émigrants sont des adultes \textbf{chômeurs} qui quittent le pays. Leur départ réduit $U$ et donc la population active ($E + U$), tandis que la population en âge de travailler diminue aussi. Le taux de participation $= (E+U)/\text{Pop.}$ diminue : par définition, tous les émigrants sont \textbf{actifs} (chômeurs), donc retirer des personnes qui sont toutes actives fait nécessairement baisser le taux de participation. Le taux de chômage baisse aussi (même mécanisme que les travailleurs découragés de la question~2). (b) est faux : le taux de participation diminue, pas augmente.
\end{reponse}

% ---------- Modèle d'offre et de demande de travail (Q9--Q16) ----------

\qcm{Dans le modèle d'offre et de demande de travail vu en cours, lequel des facteurs suivants déplace la courbe de \textbf{demande} de travail\,?}
\choix{
   \item Une hausse de la population en âge de travailler
   \item Une baisse des impôts sur le revenu des particuliers
   \item Une augmentation de la richesse des ménages
   \item Une hausse de la productivité des travailleurs
}

\begin{reponse}
\textbf{(d)} Les déterminants de la \textbf{demande} de travail vus en cours sont : le cycle économique, les impôts sur la masse salariale, le coût du capital et la \textbf{productivité de la main-d'{\oe}uvre}. Une hausse de la productivité augmente la valeur de chaque travailleur pour l'entreprise $\to$ la demande se déplace vers la droite. Les trois autres options sont des facteurs d'\textbf{offre} : (a)~la démographie, (b)~les impôts sur le revenu et (c)~la richesse des ménages déplacent la courbe d'offre de travail, pas la demande.
\end{reponse}

\qcm{Selon le modèle du marché du travail, lequel des éléments suivants devrait entraîner une \textbf{baisse de l'emploi} et une \textbf{hausse du salaire réel} à l'équilibre\,?}
\choix{
   \item Une baisse de l'offre de main-d'{\oe}uvre
   \item Une augmentation de l'offre de main-d'{\oe}uvre
   \item Une baisse de la demande de main-d'{\oe}uvre
   \item Une augmentation de la demande de main-d'{\oe}uvre
}

\begin{reponse}
\textbf{(a)} Pour que l'emploi baisse \textbf{et} le salaire augmente simultanément, il faut un déplacement de la courbe d'\textbf{offre} vers la \textbf{gauche}. Quand l'offre de travail diminue (ex. : départ à la retraite, émigration), il y a moins de travailleurs disponibles $\to$ les entreprises doivent offrir des salaires plus élevés pour attirer les travailleurs restants, mais l'emploi total diminue. (c) donnerait emploi$\downarrow$ et salaire$\downarrow$; (b) et (d) donneraient emploi$\uparrow$.
\end{reponse}

\qcm{La participation massive des femmes au marché du travail après 1950 a déplacé l'offre de travail vers la droite. À \textbf{court terme} ($K$ fixe), le modèle prédit :}
\choix{
   \item Une hausse des salaires réels et une hausse de l'emploi
   \item Aucun effet, car le capital s'ajuste immédiatement
   \item Une baisse des salaires réels et une hausse de l'emploi
   \item Une hausse des salaires réels et une baisse de l'emploi
}

\begin{reponse}
\textbf{(c)} L'offre de travail se déplace vers la \textbf{droite} (plus de travailleurs disponibles à chaque niveau de salaire). À court terme, le capital est fixe, donc la demande ne bouge pas. Le nouvel équilibre se trouve en descendant le long de la courbe de demande : l'emploi augmente ($L \uparrow$) et le salaire réel baisse ($w \downarrow$). C'est exactement le diagramme de court terme vu en cours pour la participation des femmes.
\end{reponse}

\qcm{Dans le modèle vu en cours, la distinction entre le court terme et le long terme sur le marché du travail repose principalement sur :}
\choix{
   \item Le temps que mettent les travailleurs à trouver un emploi
   \item L'intervention du gouvernement, qui n'a lieu qu'à long terme
   \item Le fait que les salaires sont rigides à court terme mais flexibles à long terme
   \item Le fait que le capital ($K$) est fixe à court terme mais s'ajuste à long terme via l'investissement
}

\begin{reponse}
\textbf{(d)} La distinction centrale est le comportement du \textbf{capital}. À court terme, le stock de capital est \textbf{fixe} : seul le travail s'ajuste, et un choc d'offre (ex. : immigration) n'affecte pas la demande de travail. À long terme, le capital \textbf{s'ajuste via l'investissement} (mécanisme de Solow) : un changement du ratio $K/L$ modifie le rendement du capital, ce qui déplace la demande de travail. (a) décrit le chômage frictionnel, un concept différent. (c) confond avec la rigidité nominale, qui est une imperfection distincte du cadre court terme / long terme.
\end{reponse}

\qcm{En récession, si les salaires nominaux sont rigides à la baisse, le résultat est :}
\choix{
   \item Une hausse des profits des entreprises
   \item Du chômage, car la quantité demandée de travail est inférieure à la quantité offerte au salaire en vigueur
   \item Une baisse du niveau des prix
   \item Une hausse de la productivité
}

\begin{reponse}
\textbf{(b)} Quand la demande de travail diminue (récession), le salaire d'équilibre devrait baisser pour rétablir le plein emploi. Mais si les salaires nominaux sont \textbf{rigides à la baisse} (conventions collectives, contrats implicites, normes sociales), le salaire reste au-dessus du nouvel équilibre. Au salaire en vigueur, la quantité demandée ($L^D$) est inférieure à la quantité offerte ($L^S$) $\to$ l'écart $L^S - L^D$ représente le \textbf{chômage involontaire}. C'est le mécanisme central illustré par le diagramme de rigidité salariale vu en cours.
\end{reponse}

\qcm{Suite à la hausse de la participation des femmes, les salaires réels avaient baissé à court terme. Que prédit le modèle à \textbf{long terme}\,?}
\choix{
   \item L'emploi est durablement plus élevé, mais les salaires reviennent à leur niveau initial grâce à l'ajustement du capital
   \item Les salaires restent définitivement plus bas, car l'offre de travail est plus grande
   \item L'emploi revient à son niveau initial, car les femmes quittent le marché
   \item Les salaires augmentent au-dessus de leur niveau initial grâce aux gains de productivité
}

\begin{reponse}
\textbf{(a)} À long terme, le capital s'ajuste. Plus de travailleurs $\to$ le ratio $K/L$ diminue $\to$ le rendement marginal du capital augmente $\to$ les entreprises investissent davantage $\to$ la demande de travail se déplace vers la \textbf{droite}. Le nouvel équilibre de long terme a un emploi \textbf{durablement plus élevé} ($L_2 > L_0$), mais les salaires \textbf{reviennent à leur niveau initial} ($w = w_0$). C'est le même mécanisme d'ajustement du capital (inspiration Solow) que pour l'immigration.
\end{reponse}

\qcm{Selon le modèle vu en cours, pourquoi l'immigration n'a-t-elle pas d'effet permanent sur les salaires à long terme\,?}
\choix{
   \item Les immigrants finissent par retourner dans leur pays d'origine
   \item Le gouvernement compense en augmentant le salaire minimum
   \item Les travailleurs natifs quittent les secteurs touchés, annulant l'effet
   \item L'investissement en capital augmente en réponse à la baisse du ratio $K/L$, ramenant la demande de travail vers la droite
}

\begin{reponse}
\textbf{(d)} C'est le mécanisme d'ajustement du capital, inspiré du modèle de Solow. À \textbf{court terme}, l'arrivée d'immigrants déplace l'offre de travail vers la droite $\to$ les salaires baissent. Mais cette baisse réduit le ratio $K/L$, ce qui augmente le \textbf{rendement marginal du capital} $\to$ les entreprises investissent davantage $\to$ $K$ augmente $\to$ la demande de travail se déplace vers la droite. À \textbf{long terme}, le capital s'est ajusté et les salaires reviennent à leur niveau initial. C'est pourquoi la plupart des études empiriques trouvent un effet proche de zéro sur les salaires à long terme.
\end{reponse}

\qcm{Les progrès dans les technologies de l'information remplacent les tâches routinières des travailleurs non qualifiés, tout en augmentant la productivité des travailleurs qualifiés. Dans le modèle d'offre et de demande de travail, cela :}
\choix{
   \item Déplace la demande vers la droite pour les qualifiés et vers la gauche pour les non-qualifiés, creusant l'écart salarial
   \item Déplace l'offre de travail non qualifié vers la gauche, car ces travailleurs se recyclent
   \item Augmente les salaires des deux groupes proportionnellement
   \item N'a aucun effet sur les inégalités salariales, car le marché s'ajuste
}

\begin{reponse}
\textbf{(a)} C'est le mécanisme du \textbf{progrès technologique biaisé en faveur des qualifications} (SBTC), illustré par les deux diagrammes côte à côte vus en cours. Les technologies de l'information augmentent la productivité des travailleurs qualifiés $\to$ leur courbe de demande se déplace vers la \textbf{droite} $\to$ leur salaire et leur emploi augmentent. En même temps, les tâches routinières des travailleurs non qualifiés sont automatisées $\to$ leur courbe de demande se déplace vers la \textbf{gauche} $\to$ leur salaire et leur emploi diminuent. L'écart salarial entre les deux groupes --- la \textit{prime de qualification} --- se creuse.
\end{reponse}

% ============================================================================
\section{Vrai, faux ou incertain}
% ============================================================================

\setcounter{shortcounter}{0}

% ---------- Indicateurs ----------

\vfi{Une baisse du taux de chômage signifie toujours que le marché du travail se porte mieux.}

\begin{reponse}
\textbf{Faux.} Le taux de chômage peut baisser pour de \guillemotleft~mauvaises raisons~\guillemotright. Par exemple, si des travailleurs découragés cessent de chercher un emploi, ils quittent la population active ($U \downarrow$, $I \uparrow$). Le taux de chômage baisse, mais aucun emploi n'a été créé. Le taux de participation diminue aussi. De même, le vieillissement fait baisser le chômage (les retraités quittent la population active) sans amélioration réelle du marché du travail. C'est pourquoi les économistes regardent les \textbf{trois} indicateurs conjointement : taux de participation, taux d'emploi \textbf{et} taux de chômage.
\end{reponse}

\vfi{Si le nombre de chômeurs et le nombre d'employés augmentent tous les deux, le taux de chômage augmente nécessairement.}

\begin{reponse}
\textbf{Faux.} Le taux de chômage $= U/(E+U)$. Si $E$ et $U$ augmentent tous les deux, l'effet sur le ratio dépend des augmentations \textbf{proportionnelles} : le taux de chômage baisse si et seulement si $\Delta U / U < \Delta E / E$ (c'est-à-dire si $E$ croît proportionnellement plus vite que $U$). Exemple : si $E$ passe de 900\,000 à 1\,000\,000 (+11.1\,\%) et $U$ de 100\,000 à 105\,000 (+5\,\%), le taux passe de 10.0\,\% à 9.5\,\%. Cela peut se produire lors d'une expansion qui crée beaucoup d'emplois tout en attirant de nouveaux chercheurs d'emploi dans la population active.
\end{reponse}

% ---------- Modèle : identification des chocs ----------

\vfi{Si la productivité des travailleurs augmente, le modèle d'offre et de demande de travail prédit que les salaires réels et l'emploi augmentent.}

\begin{reponse}
\textbf{Vrai.} Une hausse de la productivité augmente la valeur marginale de chaque travailleur pour l'entreprise $\to$ la courbe de \textbf{demande} de travail se déplace vers la \textbf{droite}. Le nouvel équilibre se trouve en montant le long de la courbe d'offre : le salaire réel augmente ($w \uparrow$) et l'emploi augmente ($L \uparrow$). Les deux variables bougent dans la \textbf{même direction} parce que c'est un déplacement de la demande : on monte le long d'une courbe d'offre à pente positive. C'est un résultat important : à long terme, c'est la croissance de la productivité qui tire les salaires vers le haut.
\end{reponse}

\vfi{Si on observe simultanément une hausse des salaires réels et une hausse de l'emploi, cela indique que l'offre de travail a augmenté.}

\begin{reponse}
\textbf{Faux.} Un déplacement de l'offre vers la \textbf{droite} donne $w \downarrow$ et $L \uparrow$ (on descend le long de la courbe de demande). Or ici les salaires \textit{augmentent}. Pour observer $w \uparrow$ et $L \uparrow$ simultanément, il faut un déplacement de la \textbf{demande} vers la droite (on monte le long de la courbe d'offre). La direction conjointe des salaires et de l'emploi permet d'identifier quelle courbe s'est déplacée : si $w$ et $L$ bougent dans le \textbf{même sens} $\to$ choc de demande; si $w$ et $L$ bougent en \textbf{sens contraire} $\to$ choc d'offre. (Ce diagnostic suppose qu'une seule courbe se déplace à la fois; si les deux bougent simultanément, l'identification est plus complexe.)
\end{reponse}

% ---------- Modèle : rigidité ----------

\vfi{Une inflation modérée (2--3\,\%) peut aider à réduire le chômage en période de récession.}

\begin{reponse}
\textbf{Vrai.} En récession, le salaire réel d'équilibre devrait baisser, mais les salaires nominaux sont rigides à la baisse (conventions collectives, résistance psychologique). Avec une inflation de 2--3\,\%, les entreprises peuvent réduire le \textbf{salaire réel} ($w/P$) sans baisser le salaire \textbf{nominal}. Exemple : si l'inflation est de 3\,\% et les salaires nominaux augmentent de 1\,\%, le salaire réel baisse de $\sim$2\,\%. Cela facilite l'ajustement du marché du travail et réduit le chômage involontaire. C'est un argument central pour une cible d'inflation \textbf{légèrement positive} (2\,\%) plutôt que 0\,\% : l'inflation \guillemotleft~lubrifie~\guillemotright\ le marché du travail.
\end{reponse}

\vfi{Un taux de chômage très bas est toujours une bonne nouvelle pour les entreprises.}

\begin{reponse}
\textbf{Incertain.} Un faible taux de chômage reflète une économie dynamique, ce qui est positif pour les revenus des entreprises (plus de consommateurs avec un emploi $\to$ plus de demande). Cependant, un marché du travail très \textbf{tendu} crée aussi des défis importants :
\begin{itemize}
   \item \textbf{Difficulté à recruter} et à retenir les employés
   \item \textbf{Pression à la hausse sur les salaires} : la main-d'{\oe}uvre est rare $\to$ les travailleurs ont un pouvoir de négociation accru
   \item \textbf{Coûts de recrutement} plus élevés et délais d'embauche plus longs
   \item Risque de mauvais \textbf{appariements} : les travailleurs n'ont pas les ressources pour chercher assez longtemps afin de trouver un bon \textit{match}
\end{itemize}
Dans le modèle, un taux de chômage très bas signifie que l'économie est proche du plein emploi $\to$ toute tentative d'expansion se heurte à une offre de travail limitée, ce qui pousse les salaires à la hausse.
\end{reponse}

\vfi{Dans le modèle vu en cours, le chômage involontaire ne peut exister que si les salaires sont rigides à la baisse.}

\begin{reponse}
\textbf{Vrai.} Dans le modèle de concurrence parfaite avec salaires \textbf{flexibles}, le marché s'ajuste toujours : le salaire monte ou descend jusqu'à ce que l'offre égale la demande ($L^S = L^D$). Il n'y a pas de chômage involontaire --- tous ceux qui veulent travailler au salaire d'équilibre trouvent un emploi. Le chômage involontaire n'apparaît que lorsque le salaire est \textbf{bloqué au-dessus de l'équilibre}, créant un excès d'offre ($L^S > L^D$). Cela se produit quand les salaires nominaux sont \textbf{rigides à la baisse} (conventions collectives, normes sociales, contrats implicites). C'est le mécanisme illustré par le diagramme avec $L^{D\prime}$ et le salaire rigide $w_0$ vu en cours.
\end{reponse}

% ---------- Modèle : applications SR/LR ----------

\vfi{Une récession réduit la demande de travail. Si les salaires sont parfaitement flexibles, le niveau d'emploi reste inchangé.}

\begin{reponse}
\textbf{Faux.} Même avec des salaires parfaitement flexibles, un déplacement de la demande de travail vers la \textbf{gauche} change l'équilibre : le salaire réel baisse ($w \downarrow$) \textbf{et} l'emploi diminue ($L \downarrow$) --- on se déplace vers le bas le long de la courbe d'offre. La flexibilité salariale empêche le \textbf{chômage involontaire} (pas d'excès d'offre, car $L^S = L^D$ au nouveau salaire), mais elle n'empêche pas la \textbf{baisse de l'emploi d'équilibre}. Certains travailleurs quittent volontairement le marché, car le salaire plus bas ne compense plus le coût de travailler. Distinction importante : la rigidité salariale \textit{aggrave} les pertes d'emploi (en ajoutant du chômage involontaire par-dessus), mais même sans rigidité, une récession fait baisser l'emploi.
\end{reponse}

\vfi{Si les salaires réels baissent et l'emploi augmente en même temps, cela est compatible avec une augmentation de l'offre de travail.}

\begin{reponse}
\textbf{Vrai.} Un déplacement de l'offre vers la \textbf{droite} (ex. : immigration, hausse de la participation) déplace l'équilibre le long de la courbe de demande à pente négative : l'emploi augmente ($L \uparrow$) et le salaire réel baisse ($w \downarrow$). C'est exactement le résultat de court terme vu en cours pour la participation des femmes et l'immigration. À l'inverse, si la demande avait augmenté, on aurait observé $w \uparrow$ \textbf{et} $L \uparrow$ (les deux dans le même sens).
\end{reponse}

\vfi{Le vieillissement de la population augmente les salaires réels de façon permanente.}

\begin{reponse}
\textbf{Faux.} Le vieillissement déplace l'offre de travail vers la \textbf{gauche} (moins de travailleurs). À \textbf{court terme}, les salaires augmentent effectivement ($w \uparrow$), car la main-d'{\oe}uvre est plus rare. Mais à \textbf{long terme}, le capital s'ajuste : moins de travailleurs $\to$ le ratio $K/L$ augmente $\to$ le rendement marginal du capital \textbf{diminue} $\to$ les entreprises investissent moins $\to$ $K$ baisse $\to$ la demande de travail se déplace vers la gauche. Le nouvel équilibre de long terme a un emploi \textbf{durablement plus faible}, mais les salaires \textbf{reviennent à leur niveau initial}. C'est le diagramme de long terme du vieillissement vu en cours.
\end{reponse}

\vfi{Si l'immigration augmente l'offre de travail en même temps qu'un boom économique augmente la demande de travail, l'effet sur l'emploi est certain mais l'effet sur les salaires est ambigu.}

\begin{reponse}
\textbf{Vrai.} Deux chocs simultanés :
\begin{itemize}
   \item \textbf{Immigration} $\to$ offre vers la droite $\to$ pression : $L \uparrow$, $w \downarrow$
   \item \textbf{Boom économique} $\to$ demande vers la droite $\to$ pression : $L \uparrow$, $w \uparrow$
\end{itemize}
Les deux chocs poussent l'emploi dans la \textbf{même direction} (hausse), donc $L$ augmente \textbf{certainement}. Mais les effets sur les salaires vont en \textbf{sens contraire} : l'immigration tire les salaires vers le bas, le boom les tire vers le haut. L'effet net sur $w$ dépend de l'ampleur relative des deux déplacements $\to$ il est \textbf{ambigu}. C'est un bon exemple de l'utilité du modèle pour raisonner sur des chocs simultanés.
\end{reponse}

\vfi{L'automatisation d'une tâche dans une entreprise réduit l'emploi total dans cette entreprise.}

\begin{reponse}
\textbf{Incertain.} Selon le cadre d'Acemoglu et Restrepo vu en cours, l'automatisation a deux effets opposés :
\begin{itemize}
   \item \textbf{Effet de déplacement :} l'automatisation remplace des travailleurs dans la tâche automatisée $\to$ emploi $\downarrow$
   \item \textbf{Effet de réintégration :} la réduction des coûts permet à l'entreprise de baisser ses prix, d'augmenter sa production et de créer de nouvelles tâches $\to$ emploi $\uparrow$
\end{itemize}
L'effet net dépend du \textbf{type de technologie} :
\begin{itemize}
   \item \textbf{Technologie transformatrice} (ex. : guichets automatiques $\to$ banques ouvrent plus de succursales $\to$ emploi total $\uparrow$) : l'effet de réintégration domine
   \item \textbf{Automatisation \guillemotleft~so-so~\guillemotright} (ex. : caisses libre-service $\to$ moins de caissiers, pas de nouvelles tâches ni de gains substantiels) : l'effet de déplacement domine
\end{itemize}
C'est la distinction vue en cours entre les technologies qui font \textbf{croître le gâteau} et celles qui ne font que redistribuer les parts.
\end{reponse}

\end{document}
